\metatitle{Non-commutative shifted symmetric functions}
\metadescription{Non-commutative shifted symmetric functions.}


\section[nonCommutativeShiftedSymmetricFunctions]{Non-commutative shifted symmetric functions}

\name{Robert Laugwitz} and \name{Vladimir Retakh} introduce
non-commutative analogues of shifted symmetric functions, such as power-sum,
complete and elementary symmetric functions, and ribbon Schur functions
\cite{LaugwitzRetakh2020}.

The ring depends on a sequence of shift parameters $a=(a_i)_{i\in\setZ}$.
The choice $a_i=0$ recovers the
\hyperref[nonComFunctions]{noncommutative symmetric functions}, and
$a_i=i-1$ gives a noncommutative analogue of the
\hyperref[schurShifted]{shifted symmetric functions} of
\name{Andrei Okounkov} and \name{Grigori Olshanski}
\cite[Ex.~2 and Ex.~5]{LaugwitzRetakh2020}.
For equally spaced parameters, the generators can be realized as rational
functions in noncommuting variables \cite[Thm.~51]{LaugwitzRetakh2020}.
For $a_i=i-1$ and commuting variables, the complete homogeneous and
elementary generators become the commutative shifted functions
$h^*_k$ and $e^*_k$ \cite[Prop.~56]{LaugwitzRetakh2020}.
As recalled in \cite[Sect.~1]{LaugwitzRetakh2020}, the commutative shifted
symmetric functions describe the center of $U(\mathfrak{gl}_n)$ and
normalized characters of symmetric groups.
For the noncommutative analogue, the authors expect applications to the
representation theory of the $0$-Hecke algebra
\cite[Sect.~1.5]{LaugwitzRetakh2020}.

\todo{Add examples showing where non-commutative shifted symmetric functions
are used, once such applications appear.  As of November 2026, OpenAlex,
Semantic Scholar and a full-text search of the local math.CO arXiv archive
found no work citing LaugwitzRetakh2020; the 0-Hecke application is only
announced in its Sect.~1.5.}
